%% Chapter 6 — Section 6.4: Discussion and Bibliography
\section{Discussion and Bibliography}
\label{sec:6.4}

The idea of using continuous-time Langevin for sampling a given target distribution can be
found in \cite{20} and \cite{98}, but it was only developed in \cite{203} where the MALA
algorithm was first proposed and analyzed. This chapter emphasized the optimization viewpoint
on Langevin Monte Carlo considered for instance in \cite{78,79,237,64}. These works establish
convergence guarantees for Langevin equation and discretizations thereof under convexity
assumptions akin to those required for convergence of gradient descent optimization \cite{34}.
The choice of step-size for MALA as well as its optimal scaling is discussed in \cite{198}.
Recent works show improvement in the scaling with dimension under appropriate convexity
assumptions \cite{139}.

As for optimization algorithms, \emph{preconditioning} Langevin Monte Carlo algorithms can
speed-up their convergence. Preconditioning improves the conditioning of the target and may be
implemented, for instance, by leveraging an ensemble of particles \cite{137,77,44}.
Relatedly, many algorithms leverage ideas from Riemannian geometry to adapt Langevin dynamics
to the geometry induced by the target distribution \cite{89,184,79}.

An important caveat of Langevin Monte Carlo algorithms is the need to evaluate gradients, which
in many applications is computationally expensive. For this reason, there has been significant
interest in studying Langevin Monte Carlo algorithms that use approximate gradients \cite{53}.
In this direction, Stochastic Gradient Langevin Dynamics (SGLD) \cite{236} provides an
important extension to ULA and MALA. Consider the setting of Bayesian inference for a
parameter $\theta$ under the assumption of i.i.d.\ data $\{y_i\}_{i=1}^K$ giving rise to a
target posterior distribution of the form
\[
  f(\theta \mid y_1, \ldots, y_K) \propto f(\theta)\prod_{i=1}^K f(y_i \mid \theta).
\]
Evaluating the log-likelihood involves computing a sum over $K$ terms, which is expensive for
large $K$. This motivates subsampling $k \ll K$ data points uniformly at random at each
iteration, which leads to the SGLD update rule
\[
  \theta^{(n+1)} = \theta^{(n)} + \epsilon_n\left\{
    \nabla\log\bigl(f(\theta^{(n)})\bigr)
    + \frac{K}{k}\sum_{i=1}^k \nabla\log\bigl(f(y_i^{(n)} \mid \theta^{(n)})\bigr)
  \right\} + \sqrt{2\epsilon_n}\,\xi^{(n)}.
\]
Here $y_i^{(n)}$ denotes the $i$-th data point of the $k$ subsampled points during the $n$-th
iteration, and $\xi^{(n)} \sim \Normal(0, I)$. Further, the step-size $\varepsilon_n$ is
chosen such that $\sum_{n=1}^\infty \epsilon_n = \infty$ and $\sum_{n=1}^\infty \epsilon_n^2
< \infty$, ensuring that we can sample from this chain \emph{without} the need for accepting
and rejecting. Note that $\epsilon_n = 1/n$ satisfies the requirements on the step size.
Another approach to design gradient-free Langevin Monte Carlo algorithms is to rely on an
ensemble of interacting particles coupled by their empirical covariance to approximate gradients
and precondition the dynamics \cite{137,77,44}.

Section~\ref{sec:6.3} provided a brief and informal introduction to Langevin equation. For
rigorous but accessible introductions to stochastic differential equations, we refer to the
textbooks \cite{177,72,153}, and for more details on Langevin equation and its use for sampling
we refer to \cite{185}. Numerical solutions are covered in \cite{125,111}. From a theoretical
perspective, several works \cite{138,191} have shown that including a non-reversible linear
drift term in Langevin equation speeds-up the convergence to the stationary distribution in
continuous time; however, the non-reversible diffusion can be more expensive to discretize
\cite{79,78}, making less obvious the algorithmic value of this idea. A general recipe for
building preconditioned non-reversible Langevin-type diffusions can be found in \cite{150},
while \cite{151} discusses Langevin methods for sampling in hidden Markov models.
